% !PRINTED 254
\weberpage{254}
Es ist aber nach (12) $\beta$ sowohl durch $e$ als auch durch $e_2$
teilbar und $\beta e e_2$ ein Quadrat; also
\[
\left(\frac{\alpha}{\beta}\right) = \left(\frac{\alpha}{ee_2}\right) = \left(\frac{\alpha}{e}\right)\left(\frac{\alpha}{e_2}\right) = \left(\frac{\delta}{e}\right)\left(\frac{\alpha}{e_2}\right) \ (\text{wegen}\ \alpha\delta \equiv 1\ \mathrm{mod}\ e);
\]
ferner ist nach (12)
\[
\left(\frac{\delta}{e}\right) = \left(\frac{c}{e}\right)\left(\frac{\partial_2}{e}\right), \quad \left(\frac{\alpha}{e_2}\right) = \left(\frac{c_2}{e_2}\right)\left(\frac{-\partial}{e_2}\right),
\]
und
\[
\left(\frac{\partial_2}{e}\right) = \left(\frac{e_2}{e}\right), \quad \left(\frac{-\partial}{e_2}\right) = \left(\frac{-e}{e_2}\right), \quad \left(\frac{e_2}{e}\right)\left(\frac{-e}{e_2}\right) = (-1)^{\frac{(e+1)(e_2-1)}{4}},
\]
also
\[
\left(\frac{\alpha}{\beta}\right) = (-1)^{\frac{(e+1)(e_2-1)}{4}}\left(\frac{c}{e}\right)\left(\frac{c_2}{e_2}\right)
\]
\[
i^{\frac{1-\beta}{2}} = i^{\frac{1-ee_2}{2}}; \quad i^{\frac{a-a_2}{2}} = i^{\frac{e-e_2}{2}}.
\]

Durch die Vertauschungen (2) geht also
\weq{14}{P_{c,\partial,a}\ \text{in}\ e^{\frac{2\pi i}{3}\frac{\lambda-1}{8}}P_{c_1,\partial,a}\ \text{und in}\ e^{\frac{2\pi i}{3}\frac{\beta(\alpha+\delta)-(\beta^2-1)\alpha\gamma}{8}}P_{c_2,\partial_2,a_2}\ \text{über}.}

Ist $n$ noch durch 3 unteilbar, so nehme man $c$ durch 3 teilbar an,
wodurch
\[
\lambda \equiv 1, \quad \alpha \equiv \delta \equiv 0\ (\mathrm{mod}\ 3)
\]
werden und die in (14) vorkommenden dritten Einheitswurzeln den Wert
1 erhalten.

Hieraus ergeben sich die Sätze:

5. \gesperrt{Ist $n$ eine ungerade Quadratzahl, $c$ durch 8 teilbar,
so sind die Größen}
\[
P^3_{c,\partial,a},
\]
\gesperrt{und ist $n$ eine durch 3 nicht teilbare ungerade
Quadratzahl, $c$ durch 24 teilbar, so sind die Größen}
\[
P_{c,\partial,a}
\]
\gesperrt{Wurzeln von Multiplikatorgleichungen erster Stufe.}

Die ersten Beispiele sind $n=9$, $n=25$.

Zur Bildung dieser Gleichungen kann man auf verschiedene Arten
gelangen. Wir wollen hier [nach \gesperrt{Kiepert}$^1$)] den Weg
gehen, daß wir nach § 69 die Wurzeln einer zum Grade $p$ ge-
%CONT>>

\begingroup
\renewcommand{\thefootnote}{1)}
\footnotetext{Zur Transformationstheorie der elliptischen
Funktionen. Crelles Journal, Bd. 87, 88, 95.}
\endgroup
